\documentclass[10pt,a4paper]{amsart}
\usepackage[margin=19mm]{geometry}
\usepackage{amsmath,amssymb,mathtools}
\usepackage[scr=rsfs,cal=euler]{mathalfa}
\usepackage{xfrac}
\usepackage[unicode,hidelinks,bookmarks=false]{hyperref}
\newcommand{\ie}{i.e.\ }
\newcommand{\cf}{cf.\ }
\newcommand{\resp}{respectively\ }
\newcommand{\Subsection}{\subsection}
\providecommand{\nobreakdash}{\nobreak\hbox{-}}
\setlength{\emergencystretch}{2em}
\multlinegap0pt
\usepackage{bm}
\usepackage{bbm}
\usepackage{dsfont}
\usepackage{xspace}
\usepackage{cancel}
\usepackage{mathtools}
\usepackage{amsthm}
\makeatletter
\@ifundefined{lemma}{\newtheorem{lemma}{Lemma}}{}
\makeatother
\newcommand{\N}{\mathbb{N}}
\title[Well-posedness and numerical analysis of an elapsed time mod]{Well-posedness and numerical analysis of an elapsed time model with strongly coupled neural networks}
\author{Mauricio Sepulveda, Nicolas Torres, Luis Miguel Villada}
\date{Source: arXiv:2310.02068v4 (CC BY 4.0)}
\begin{document}
\maketitle

\noindent\emph{Source and licence.} Well-posedness and numerical analysis of an elapsed time model with strongly coupled neural networks. Version arXiv:2310.02068v4 (CC BY 4.0), distributed under the Creative Commons Attribution 4.0 International licence, as recorded in the arXiv licence metadata for this exact version and shown on the abstract page https://arxiv.org/abs/2310.02068v4. Full rights notice: licenses/arxiv-2310.02068-CC-BY-4.0.txt.\par\smallskip

\noindent\textit{Editorial scope.} Counted unit: the Lemma whose source label is \texttt{BV-lem-d} in the article (source0.tex lines 987--995, source label \texttt{BV-lem-d}), delivered together with the article's own complete original proof of it (source0.tex lines 996--1092, one \texttt{proof} environment of 97 lines). No mathematics is written, repaired or re-derived: statement and proof are the article's own text line for line. Required context retained uncounted: CFL (lines 418--420); eq:findX (lines 888--890). Every definition, display and label of those lines is retained unchanged. Omitted from this package and not redistributed: 1--417, 421--887, 891--986, 1093--1332. Nothing inside a retained excerpt is omitted or altered: every retained body is delivered exactly as the article's lines are written, so no omission receipt of any kind accompanies this package. Citations: the retained excerpts carry no citation command at all, so no bibliography entry of the article is transcribed and no reference is redistributed. Reference adaptation: none was needed and none was made; every reference inside the delivered unit resolves inside it, so no unresolved reference or citation is produced. Printed slips kept literal: no printed word, symbol, hypothesis or number of the counted statement or of its proof was altered. Layout and class: the article's own class is \texttt{article} with packages including a4wide, graphicx, amssymb, amsmath, amsthm, mathrsfs, xfrac, mathtools, enumerate, color, hyperref, array, amsfonts, psfrag; the wrapper is the lane's pinned \texttt{amsart} editorial wrapper at 10pt on A4, which restates the theorem-like environments the retained text needs and adds the article's own macro definitions that the retained text needs (\texttt{\textbackslash N}), with extra wrapper packages (bm, bbm, dsfont, xspace, cancel, mathtools). The delivered numbering is the unit's own. Assets: no retained span contains a figure environment, a graphics-inclusion command, a PDF-page inclusion, a table or a TikZ picture; no image file is copied into this package, the package declares no asset dependency and no image file is redistributed with it. Licence: version arXiv:2310.02068v4 of this article is distributed under the Creative Commons Attribution 4.0 International licence, as recorded in the arXiv licence metadata for this exact version and shown on the abstract page https://arxiv.org/abs/2310.02068v4. A full rights notice is delivered at licenses/arxiv-2310.02068-CC-BY-4.0.txt. Characters: every retained excerpt is pure ASCII, so the delivered engine, XeLaTeX with its default Latin Modern text font, reports no missing character.
 \begin{equation}\label{CFL}
     \Delta t \le \left(\frac{1}{\Delta s}+\|p\|_\infty\right)^{-1}.
 \end{equation}

\begin{equation}\label{eq:findX}
    X^m=\frac{\Delta t}{2} \left( \Delta s \sum_{j\in\N}p(s_j,X^m)n_j^m \alpha_0 + \sum_{k=0}^{m-1} N^k\alpha_{m-k} \right),
\end{equation}

\begin{lemma}{(\textbf{BV}-estimate)}
\label{BV-lem-d}
Assume that $TV(n^0)<\infty$ and the CFL condition \eqref{CFL}. Then there exist constants $C_1,C_2>0$ (depending only $p,\alpha$ and the norms of $n^0$) such that for $m\in\N$ we have
\begin{equation}
    TV(n^m)\le e^{C_1T} TV(n^0)+C_2(e^{C_1T}-1),
    \label{BV_estimate-d}
\end{equation}
with  $T=m\Delta t$ and $TV(n^m)= \displaystyle\sum_{j=0}^\infty|n_{j+1}^{m}-n_{j}^{m}|$.
\end{lemma}

\begin{proof}
using notation $\Delta^{+}n_j^m=n_{j+1}^{m}-n_{j}^{m}$, we have
\begin{equation*}
    \begin{split}
      \Delta^{+}n_j^{m+1}&=  \Delta^{+}n_j^{m}-\frac{\Delta t}{\Delta s}\left(\Delta^{+}n_j^{m}-\Delta^{+}n_{j-1}^{m}\right)-\Delta t p(s_{j+1},X^m)\Delta^{+}n_j^{m}
       -\Delta t n_j^m(p(s_{j+1},X^m)-p(s_{j},X^m)).
    \end{split}
\end{equation*}
Next, from the CFL condition \eqref{CFL} we obtain
\begin{equation*}
     |\Delta^{+}n_j^{m+1}| \le \left(1-\Delta t\left( \frac{1}{\Delta s}+p(s_{j+1},X^m) \right) \right)|\Delta^{+}n_j^{m}|
     +\frac{\Delta t}{\Delta s}|\Delta^{+}n_{j-1}^{m}|+\Delta t\Delta s\|\partial_s p\|_\infty n_j^m.
\end{equation*}
By summing over all $j\ge 1$ we deduce that
\begin{equation}
\label{Delta+nm}
\sum_{j=1}^\infty|\Delta^{+}n_j^{m+1}|\le
\sum_{j=1}^\infty|\Delta^{+}n_j^{m}|
+ \frac{\Delta t}{\Delta s} |\Delta^{+}n_{0}^{m}|+\Delta t\|\partial_s p\|_\infty\|n^0\|_1. 
\end{equation}
We now take into account the boundary term. From the numerical scheme we have
\begin{equation*}
\label{Delta+n0-d}
\begin{split}
     |\Delta^{+}n_0^{m+1}| & = |n_1^{m+1}-N^{m+1}|\\
     & \le \left(1-\frac{\Delta t}{\Delta s}\right)|n_1^m-N^m|+  |N^{m+1}- N^m|+\Delta t p(s_1,N^m) n^m_1\\
     &\le \left(1-\frac{\Delta t}{\Delta s}\right)|\Delta^+ n_0^m|+  |N^{m+1}-N^m|+\Delta t \|p\|_\infty \|n^0\|_\infty.
\end{split} 
\end{equation*}

Next we estimate the second term of the last inequality. First note that
\begin{equation*}
    \begin{split}
      N^{m+1}- N^{m}& =\Delta s \sum_{j\in\N}p(s_j,X^{m+1})n_j^{m+1}-\Delta s \sum_{j\in\N} p(s_j,X^{m})n_j^{m}\\
    &=\Delta s \sum_{j\in\N}(p(s_j,X^{m+1})-p(s_j,X^{m}))n_j^{m+1}+\Delta s \sum_{j\in\N}p(s_j,X^m)(n_j^{m+1}-n_j^m),
    \end{split}
\end{equation*}
thus we have
\begin{equation}
\label{N1-N2-d}
    \begin{split}
      |N^{m+1}- N^{m}|&\le \|\partial_X p\|_\infty \|n^0\|_1 |X^{m+1}-X^{m}|+\|p\|_\infty\sum_{j\in\N}\Delta s|n^{m+1}_j-n^m_j|\\
      &\le  \|\partial_X p\|_\infty \|n^0\|_1 |X^{m+1}-X^{m}| + \|p\|_\infty\Delta t\left(  \sum_{j\in\N} | n_j^{m}-n_{j-1}^m|  + \|p\|_\infty \sum_{j\in\N}  \Delta s n_j^m \right)\\
      &\le  \|\partial_X p\|_\infty \|n^0\|_1 |X^{m+1}-X^{m}| + \|p\|_\infty\Delta t  \sum_{j\in\N} |\Delta^+ n^m_{j-1}| + \Delta t\|p\|_\infty^2\|n^0\|_1.
    \end{split}
\end{equation}
%\begin{eqnarray}
%  |N^{m+1}- N^{m}|&\le&\Delta s \sum_{j\in\N} left\{|p_X(s_j,\vartheta)||X^{m+1}-X^m|n_j^{m+1} + p(s_j,X^m)|n_j^{m+1}-n_j^m|  \right\}  \\
%  &\le&\|p_X\|_{\infty}\|n^{m+1}\|_1|X^{m+1}-X^m|+\Delta s \|p\|_{\infty} \sum_{j=1}^\infty|nj^{m+1}-n_j^m|.
%\end{eqnarray}
Similarly for $|X^{m+1}-X^{m}|$ we get by using Equation \eqref{eq:findX}
\begin{equation*}
    \begin{split}
       |X^{m+1}-X^m|&\le \frac{\Delta t}{2}\alpha_0\|\partial_X p\|_\infty\|n^0\|_1 |X^{m+1}-X^{m}|+\frac{\Delta t}{2}\alpha_0\|p\|_\infty\sum_{j\in\N}\Delta s|n^{m+1}_j-n^m_j|\\
       &\qquad +\frac{\Delta t}{2}\left|\sum_{k=0}^m N^k\alpha_{m+1-k}-\sum_{k=0}^{m-1} N^k\alpha_{m-k}\right|,
    \end{split}
\end{equation*}
so we get
\begin{equation*}
    \begin{split}
       |X^{m+1}-X^m|&\le \frac{\Delta t}{2}\alpha_0\left(\|\partial_X p\|_\infty\|n^0\|_1 |X^{m+1}-X^{m}|+\Delta t\|p\|_\infty\sum_{j\in\N} |\Delta^+ n^m_{j-1}| +\Delta t\|p\|_\infty^2\|n^0\|_1\right)\\
       &\qquad +\frac{\Delta t}{2} N^m\alpha_1+\frac{\Delta t}{2}\|p\|_\infty\|n^0\|_1\sum_{k=0}^{m}|\alpha_{m+1-k}-\alpha_{m-k}|,
    \end{split}
\end{equation*}
and therefore we have the following estimate
\begin{equation}
\begin{split} \label{X1-X2-d}
     |X^{m+1}-X^m| &\le \frac{\Delta t}{2(1-\frac{\Delta t}{2}\alpha_0\|\partial_X p\|_\infty\|n^0\|_1)}\left(\Delta t\alpha_0\|p\|_\infty\sum_{j\in\N} |\Delta^+ n^m_{j-1}| +\Delta t\alpha_0\|p\|_\infty^2\|n^0\|_1 \right. \\
     &\left.\phantom{\sum_{j\in\N}}+
    \|p\|_\infty\|n^0\|_1\|\alpha\|_\infty+\|p\|_\infty\|n^0\|_1 TV(\alpha)\right).
\end{split}
\end{equation}
By plugging \eqref{X1-X2-d} into \eqref{N1-N2-d} we obtain for $\Delta t$ small
\begin{equation*}
|N^{m+1}-N^m|\le A_1\Delta t\sum_{j\in\N}|\Delta^+ n^m_{j-1}|+B_1\Delta t,
\end{equation*}
where $A_1,B_1>0$ are constants depending on $p,\alpha$ and the norms of $n^0$. So that the boundary term is estimated as follows
\begin{equation}
\label{Delta+n0}
\begin{split}
     |\Delta^{+}n_0^{m+1}|
     &\le \left(1-\frac{\Delta t}{\Delta s}\right)|\Delta^+ n_0^m|+  A_1\Delta t\sum_{j\in\N}|\Delta^+ n^m_{j-1}|+B_2\Delta t,
\end{split}
\end{equation}
and by adding \eqref{Delta+n0} with \eqref{Delta+nm}, we obtain
\begin{equation}
\sum_{j=0}^\infty|\Delta^{+}n_j^{m+1}|\le
(1+C_1\Delta t)\sum_{j=0}^\infty|\Delta^{+}n_j^{m}|+C_2\Delta t,
\end{equation}
where $C_1,C_2>0$ constants depending only on $p,\alpha$ and the norms of $n^0$.

Finally, proceeding recursively on $m$, we obtain
\begin{equation*}
 \sum_{j=0}^\infty|\Delta^{+}n_j^m|\le (1+C_1\Delta t)^{m}\sum_{j=0}^\infty|\Delta^{+}n_j^{0}|+\frac{C_2}{C_1}\left((1+C_1\Delta t)^{m}-1\right),
\end{equation*}
and the estimate \eqref{BV_estimate-d} readily follows.
\end{proof}

\end{document}
